Nach Archimedes ist die Auftriebskraft gleich der Gewichtskraft des verdrängten Wassers ($\rho_\text{W} = \rhoWO$, FoTaBe S.~187). Die Kugel verdrängt ihr eigenes Volumen
\begin{align}
 V &= \VKF = \frac{4}{3}\pi\cdot\left(\rK\right)^3 = \VK
\end{align}
Damit ist
\begin{align}
 F_\text{A} &= \FAF \\
   &= \rhoW\cdot\gE\cdot\VK \\
   &= \FA \approx \result{\FAP{0}{2}}
\end{align}
Das Material der Kugel spielt keine Rolle; eine gleich grosse Kugel aus Holz, die man unter Wasser hält, erfährt denselben Auftrieb.
